\section{Scope and interpretation boundary}

This Supplementary Information is the evidence-routing guide for the extended Taxila manuscript. It identifies the complete records that support the condensed methods, results and displays in the main paper. The analytical outputs describe relative landscape pressure, terrain and hydrological susceptibility, regional climate forcing and relative field-inspection priority. They do not demonstrate monument damage, condition, structural safety, annual loss probability or a legally defined risk zone. The 250, 500 and 1,000~m circles are analytical supports rather than official property or buffer boundaries. WorldCover-derived classes are proxy land-cover labels and are not independent heritage-condition observations.

The final validated experiment package is authoritative. Legacy summaries and remembered values must not be used when they differ from the claim-evidence matrix, discrepancy record or final tables. Saraikala remains in the documentary inventory but is excluded from coordinate-dependent summaries because its geometry is unresolved.

\section{Final display allocation}

The main manuscript uses 18 non-duplicative display items: 14 figures and four tables. The Supplementary Information uses four additional figures and seven detailed tables. Source identifiers below refer to the frozen experiment package. Main figures retain only the scientific argument needed for interpretation; complete component, scale, climate, sensitivity and model records remain supplementary.

\subsection{Main-manuscript display sequence}

\setcounter{table}{-2}

\begin{longtable}{p{0.10\textwidth}p{0.31\textwidth}p{0.48\textwidth}}
\toprule
ID & Display & Verified source or construction basis \\
\midrule
\endfirsthead
\toprule
ID & Display & Verified source or construction basis \\
\midrule
\endhead
F1 & Study area, component inventory and analytical supports & \path{12_maps/figure_01_study_area_context_map.pdf} \\
F2 & CHIP workflow and inference boundaries & Run sequence, configurations and \path{16_validation/experiment_status_matrix.csv}; contains no new result \\
F3 & Five-epoch relative landscape-pressure surfaces & \path{12_maps/figure_02_five_epoch_landscape_pressure_maps.pdf} \\
F4 & Endpoint adverse-spectral convergence & \path{12_maps/figure_03_spectral_pressure_convergence.pdf} \\
F5 & Annual climate trends and moving-block uncertainty & \path{11_figures/figure_05_annual_climate_trends.pdf} \\
F6 & Landsat-epoch antecedent weather & \path{11_figures/figure_07_epoch_specific_weather_forcing.pdf} \\
F7 & Open-Meteo and NASA POWER comparison & \path{11_figures/figure_09_open_meteo_nasa_power_comparison.pdf} \\
F8 & Terrain, flow and drainage susceptibility & \path{12_maps/figure_11_terrain_flow_drainage_susceptibility.pdf} \\
F9 & E2024 integrated 500~m inspection priority & \path{12_maps/figure_12_integrated_component_priority_map.pdf} \\
F10 & Component ranks across analytical supports & \path{11_figures/figure_13_component_rankings_by_scale.pdf} \\
F11 & Buffered proxy-model performance and calibration & \path{11_figures/figure_14_proxy_model_performance_calibration.pdf} \\
F12 & Indicator redundancy & \path{11_figures/figure_15_indicator_redundancy_matrix.pdf} \\
F13 & Weight and spatial-block rank uncertainty & \path{11_figures/figure_16_rank_stability_uncertainty.pdf} \\
F14 & Factor and domain ablation & \path{11_figures/figure_17_factor_domain_ablation.pdf} \\
T1 & Data sources, periods, supports and roles & Source manifest and reproducibility data dictionary \\
T2 & Experiment and evidence-completeness status & Experiment-status and unsuccessful-analysis records \\
T3 & Indicator transformations, formulae and weights & \path{13_tables/composite_formulas_weights_normalisation.csv} \\
T4 & Headline findings and inference boundaries & Final claim-evidence matrix and cited numerical outputs \\
\bottomrule
\end{longtable}

\subsection{Supplementary display sequence}

\begin{longtable}{p{0.10\textwidth}p{0.31\textwidth}p{0.48\textwidth}}
\toprule
ID & Display & Verified source output \\
\midrule
\endfirsthead
\toprule
ID & Display & Verified source output \\
\midrule
\endhead
F S1 & Monthly Open-Meteo climatology & \path{11_figures/figure_04_open_meteo_climatology.pdf} \\
F S2 & Seasonal and monsoon climate anomalies & \path{11_figures/figure_06_seasonal_monsoon_anomalies.pdf} \\
F S3 & Climate-extreme timeline, 1991--2025 & \path{11_figures/figure_08_climate_extremes_timeline.pdf} \\
F S4 & Exploratory weather-spectral lag associations & \path{11_figures/figure_10_weather_spectral_lag_associations.pdf} \\
T S1 & Component inventory and geometry status & \path{06_processed_vectors/taxila_components_wgs84.geojson} \\
T S2 & Complete ranks and scores at 250, 500 and 1,000~m & \path{13_tables/component_rankings_250m_500m_1000m.csv} \\
T S3 & Climate trends and baseline sensitivity & \path{14_statistics/open_meteo_climate_trends_block_bootstrap.csv}; \path{14_statistics/climate_baseline_period_sensitivity.csv} \\
T S4 & Complete weather-product comparison & \path{13_tables/open_meteo_nasa_power_comparison.csv} and product-agreement tables \\
T S5 & Complete factor and domain ablation & \path{13_tables/complete_factor_domain_ablation.csv}; \path{13_tables/component_factor_domain_ablation_ranks.csv} \\
T S6 & Weight and spatial-block rank distributions & Monte Carlo rank table and \path{14_statistics/component_rank_spatial_block_bootstrap.csv} \\
T S7 & Proxy-model metrics, calibration and confusion records & Buffered comparison, calibration, class distribution and primary confusion matrix tables \\
\bottomrule
\end{longtable}

\section{S1: Spatial inventory and analytical supports}

\setcounter{table}{0}

Table~S1 must report all 18 documentary component records, official identifiers, geometry status and coordinate-dependent inclusion. It must state that 17 components were geolocated, that Saraikala was not ranked, and that point-centred analytical circles are not official polygons. Figure~1 in the main paper should carry the cartographic context, coordinate reference system, scale bar, north arrow, component names and transparent NoData treatment.

\begingroup
\scriptsize
\setlength{\tabcolsep}{3pt}
\begin{longtable}{p{0.10\textwidth}p{0.30\textwidth}p{0.24\textwidth}rrp{0.08\textwidth}}
\caption{Complete Taxila component inventory and analytical geometry status.}\label{tab:s-component-inventory}\\
\toprule
ID & Component & Coordinate status & Longitude & Latitude & Included \\
\midrule
\endfirsthead
\caption[]{Complete Taxila component inventory and analytical geometry status. (continued)}\\
\toprule
ID & Component & Coordinate status & Longitude & Latitude & Included \\
\midrule
\endhead
\midrule
\multicolumn{6}{r}{Continued on next page}\\
\endfoot
\bottomrule
\endlastfoot
139-001 & Khanpur Cave & official\_web\_point & 72.909294 & 33.799200 & Yes \\
139-002 & Saraikala, prehistoric mound & official\_id\_geometry\_missing &  &  & No \\
139-003 & Bhir Mound & official\_web\_point & 72.819331 & 33.743944 & Yes \\
139-004 & Sirkap (fortified city) & official\_web\_point & 72.829878 & 33.753403 & Yes \\
139-005 & Sirsukh (fortified ruined city) & official\_web\_point & 72.849739 & 33.772256 & Yes \\
139-006 & Dharmarajika stupa and monastery & official\_web\_point & 72.842667 & 33.744669 & Yes \\
139-007 & Khader Mohra (Akhuri) & official\_web\_point & 72.860722 & 33.760806 & Yes \\
139-008 & Kalawan group of buildings & official\_web\_point & 72.839753 & 33.746292 & Yes \\
139-009 & Giri complex of monuments & official\_web\_point & 72.883325 & 33.729089 & Yes \\
139-010 & Kunala stupa and monastery & official\_web\_point & 72.831444 & 33.750553 & Yes \\
139-011 & Jandial complex & official\_web\_point & 72.828914 & 33.764456 & Yes \\
139-012 & Lalchak and Badalpur Buddhist stuppa & official\_web\_point & 72.868019 & 33.781803 & Yes \\
139-013 & Mohra Moradu stupa and monastery & official\_web\_point & 72.861053 & 33.761097 & Yes \\
139-014 & Pippala stupa and monastery & official\_web\_point & 72.866686 & 33.765922 & Yes \\
139-015 & Jaulian stupa and monastery & official\_web\_point & 72.875311 & 33.765042 & Yes \\
139-016 & Lalchak mounds & official\_web\_point & 72.849469 & 33.774450 & Yes \\
139-017 & Buddhist remains around Bhallar stupa & official\_web\_point & 72.825358 & 33.813497 & Yes \\
139-018 & Giri Mosque and tombs & official\_web\_point & 72.880522 & 33.727769 & Yes \\
\end{longtable}
\noindent\textit{Note:} Coordinates are WGS 84 point locations transcribed from the official component map. Saraikala remains in the inventory but is excluded from coordinate-dependent analysis. Analytical circles are not official property polygons or legal buffers.
\endgroup


\begingroup
\scriptsize
\setlength{\tabcolsep}{3pt}
\begin{longtable}{p{0.09\textwidth}p{0.29\textwidth}rrrrrr}
\caption{Complete component ranks and relative scores at 250, 500 and 1,000 m analytical supports.}\label{tab:s-component-ranks}\\
\toprule
ID & Component & R250 & S250 & R500 & S500 & R1000 & S1000 \\
\midrule
\endfirsthead
\caption[]{Complete component ranks and relative scores at 250, 500 and 1,000 m analytical supports. (continued)}\\
\toprule
ID & Component & R250 & S250 & R500 & S500 & R1000 & S1000 \\
\midrule
\endhead
\midrule
\multicolumn{8}{r}{Continued on next page}\\
\endfoot
\bottomrule
\endlastfoot
139-001 & Khanpur Cave & 15.0 & 0.400 & 14.0 & 0.431 & 11.0 & 0.467 \\
139-003 & Bhir Mound & 14.0 & 0.412 & 17.0 & 0.340 & 17.0 & 0.308 \\
139-004 & Sirkap (fortified city) & 10.0 & 0.474 & 12.0 & 0.480 & 15.0 & 0.395 \\
139-005 & Sirsukh (fortified ruined city) & 8.0 & 0.479 & 9.0 & 0.501 & 10.0 & 0.485 \\
139-006 & Dharmarajika stupa and monastery & 3.0 & 0.521 & 5.0 & 0.536 & 6.0 & 0.539 \\
139-007 & Khader Mohra (Akhuri) & 7.0 & 0.482 & 6.0 & 0.535 & 3.0 & 0.575 \\
139-008 & Kalawan group of buildings & 12.0 & 0.453 & 7.0 & 0.519 & 9.0 & 0.513 \\
139-009 & Giri complex of monuments & 4.0 & 0.508 & 1.0 & 0.591 & 2.0 & 0.578 \\
139-010 & Kunala stupa and monastery & 5.0 & 0.494 & 10.5 & 0.501 & 13.0 & 0.441 \\
139-011 & Jandial complex & 16.0 & 0.357 & 16.0 & 0.349 & 16.0 & 0.394 \\
139-012 & Lalchak and Badalpur Buddhist stuppa & 17.0 & 0.319 & 15.0 & 0.392 & 12.0 & 0.455 \\
139-013 & Mohra Moradu stupa and monastery & 11.0 & 0.464 & 4.0 & 0.539 & 5.0 & 0.543 \\
139-014 & Pippala stupa and monastery & 6.0 & 0.489 & 10.5 & 0.501 & 8.0 & 0.537 \\
139-015 & Jaulian stupa and monastery & 2.0 & 0.567 & 3.0 & 0.568 & 1.0 & 0.593 \\
139-016 & Lalchak mounds & 13.0 & 0.438 & 13.0 & 0.446 & 14.0 & 0.429 \\
139-017 & Buddhist remains around Bhallar stupa & 9.0 & 0.478 & 8.0 & 0.517 & 7.0 & 0.537 \\
139-018 & Giri Mosque and tombs & 1.0 & 0.622 & 2.0 & 0.573 & 4.0 & 0.569 \\
\end{longtable}
\noindent\textit{Note:} R denotes rank and S denotes relative field-inspection-priority score. Scores are conditional on the declared hierarchy and component set.
\endgroup


\section{S2: Landsat record, quality control and common support}

This section must preserve the complete 55-scene inventory, sensor allocation, three-year post-monsoon windows, quality-flag exclusions, Collection~2 scaling, valid-observation counts and common-support construction. Landsat surface-reflectance provenance and cross-sensor interpretation are anchored by established processing and product studies \citep{masek2006ledaps,vermote2016landsat,roy2016landsat}. The workflow used Collection~2 quality flags; Fmask literature documents the broader importance of cloud and shadow control but does not describe an extra algorithm executed in this study \citep{zhu2012fmask,zhu2015fmask}. Denser breakpoint and trajectory methods provide methodological context for what a five-epoch design cannot resolve \citep{verbesselt2010bfast,kennedy2010landtrendr}.

\section{S3: Spectral transformations and convergence}

This section must give the executed equations, band mapping, pressure direction, normalisation, percentile thresholds and common-support denominator. It must preserve exact endpoint thresholds and the counts for zero, one, two and three convergent criteria. The full component-by-epoch summaries belong here. Captions must call the outputs landscape pressure or spectral convergence, never condition, damage or risk.

\section{S4: Climate provenance, trends and product sensitivity}

Figures~S1--S4 and Tables~S3--S4 provide the complete monthly, seasonal, annual, extreme, baseline-sensitivity, product-comparison and acquisition-window records. ERA5 and ERA5-Land publications establish the reanalysis context but do not convert a gridded product into local station truth \citep{hersbach2020era5,munozsabater2021era5land}. Trend interpretation must retain the consequences of serial dependence and multiple comparisons \citep{hamed1998mk,yue2002autocorrelation,wilks2016stippling}. The moving-block design is a dependence-preserving resampling choice, not an assertion that five years is the uniquely correct block length \citep{politis1994stationary}. All five-epoch weather-spectral correlations remain exploratory, without causal language or formal significance claims.

\begin{figure}[p]
\centering
\includegraphics[width=0.94\textwidth]{figure_s01_climatology.pdf}
\caption{Monthly Open-Meteo climatology for the four weather cells associated with the 17 mapped Taxila components. The 1991 to 2020 baseline summarises the seasonal cycle of precipitation, temperature, humidity, soil moisture and reference evapotranspiration using each variable's native units. Values describe regional gridded forcing, not station observations or 30 m weather.}
\label{fig:s-climatology}
\end{figure}

\begin{figure}[p]
\centering
\includegraphics[width=0.90\textwidth]{figure_s02_seasonal_anomalies.pdf}
\caption{Seasonal and monsoon climate anomalies from the Open-Meteo 1991 to 2025 record. Anomalies are referenced to the matching 1991 to 2020 climatological aggregation. Precipitation is in mm and temperature in degrees C. The panels show interannual variability and do not constitute a future projection or component-scale hazard record.}
\label{fig:s-seasonal}
\end{figure}

\begin{figure}[p]
\centering
\includegraphics[width=0.91\textwidth]{figure_s03_extreme_timeline.pdf}
\caption{Timeline of selected regional climate-extreme metrics, 1991 to 2025. Metrics were calculated from the complete Open-Meteo daily series after quality control. Counts and magnitudes follow the definitions in the reproducibility data dictionary. The figure describes gridded historical variability and does not establish event occurrence or impact at a heritage component.}
\label{fig:s-extremes}
\end{figure}

\begin{figure}[p]
\centering
\includegraphics[width=0.96\textwidth]{figure_s04_weather_spectral_associations.pdf}
\caption{Exploratory associations between antecedent weather summaries and five satellite-composite spectral summaries. Coefficients use only five independent epochs; leave-one-epoch ranges show sensitivity to omitting one composite. The display is descriptive, contains no confirmatory significance claim and does not attribute spectral change to climate forcing.}
\label{fig:s-weather-spectral}
\end{figure}

\FloatBarrier
\begingroup
\scriptsize
\setlength{\tabcolsep}{3pt}
\begin{longtable}{p{0.18\textwidth}p{0.11\textwidth}p{0.25\textwidth}r r p{0.18\textwidth}r}
\caption{Complete Open-Meteo trend estimates and baseline-period sensitivity.}\label{tab:s-climate-trends}\\
\toprule
Analysis & Period & Metric & $n$ & Slope yr$^{-1}$ & Block 95\% interval & $p$ \\
\midrule
\endfirsthead
\caption[]{Complete Open-Meteo trend estimates and baseline-period sensitivity. (continued)}\\
\toprule
Analysis & Period & Metric & $n$ & Slope yr$^{-1}$ & Block 95\% interval & $p$ \\
\midrule
\endhead
\midrule
\multicolumn{7}{r}{Continued on next page}\\
\endfoot
\bottomrule
\endlastfoot
Full 1991--2025 & annual & precipitation sum mm & 35 & -10.3037 & [-17.9150, -2.6679] & 0.0360 \\
Full 1991--2025 & annual & rx1day mm & 35 & -0.6233 & [-1.5426, 0.4034] & 0.2711 \\
Full 1991--2025 & annual & rx5day mm & 35 & -1.7354 & [-3.3891, -0.2019] & 0.1105 \\
Full 1991--2025 & annual & heavy precipitation days & 35 & -0.1429 & [-0.2266, -0.0582] & 0.0140 \\
Full 1991--2025 & annual & consecutive dry days & 35 & 0.1667 & [-0.1417, 0.4322] & 0.3189 \\
Full 1991--2025 & annual & temperature mean c & 35 & 0.0418 & [0.0099, 0.0731] & 0.0350 \\
Full 1991--2025 & annual & annual or period max tmax c & 35 & 0.0333 & [-0.0079, 0.0778] & 0.2594 \\
Full 1991--2025 & annual & hot days & 35 & 0.1818 & [-0.1728, 0.5227] & 0.4085 \\
Full 1991--2025 & annual & heatwave max days & 35 & 0.0323 & [-0.0449, 0.1144] & 0.4326 \\
Full 1991--2025 & annual & relative humidity mean percent & 35 & -0.1311 & [-0.2852, 0.0255] & 0.2594 \\
Full 1991--2025 & annual & soil moisture mean m3m3 & 35 & -0.0013 & [-0.0022, -0.0003] & 0.0350 \\
Full 1991--2025 & annual & et0 sum mm & 35 & 2.2850 & [-1.2405, 5.5941] & 0.3189 \\
Full 1991--2025 & DJF & precipitation sum mm & 34 & -2.5913 & [-5.7242, 0.2144] & 0.2086 \\
Full 1991--2025 & DJF & temperature mean c & 34 & 0.0334 & [0.0100, 0.0584] & 0.0350 \\
Full 1991--2025 & DJF & rx1day mm & 34 & -0.3133 & [-0.8050, 0.1852] & 0.3189 \\
Full 1991--2025 & DJF & consecutive dry days & 34 & 0.0667 & [-0.1800, 0.3084] & 0.4948 \\
Full 1991--2025 & MAM & precipitation sum mm & 35 & -3.7432 & [-8.8320, 1.2465] & 0.2711 \\
Full 1991--2025 & MAM & temperature mean c & 35 & 0.0463 & [-0.0140, 0.1095] & 0.2711 \\
Full 1991--2025 & MAM & rx1day mm & 35 & -0.7421 & [-1.0366, -0.4565] & 0.0070 \\
Full 1991--2025 & MAM & consecutive dry days & 35 & 0.0000 & [-0.2105, 0.2308] & 1.0000 \\
Full 1991--2025 & Monsoon & precipitation sum mm & 35 & -2.5130 & [-7.1320, 2.0571] & 0.3911 \\
Full 1991--2025 & Monsoon & temperature mean c & 35 & 0.0483 & [0.0345, 0.0640] & 0.0070 \\
Full 1991--2025 & Monsoon & rx1day mm & 35 & -0.4023 & [-1.7419, 0.9624] & 0.4948 \\
Full 1991--2025 & Monsoon & consecutive dry days & 35 & -0.0294 & [-0.1095, 0.0554] & 0.4326 \\
Full 1991--2025 & Post-monsoon & precipitation sum mm & 35 & -0.2773 & [-1.4308, 0.8196] & 0.5027 \\
Full 1991--2025 & Post-monsoon & temperature mean c & 35 & 0.0417 & [0.0057, 0.0796] & 0.0886 \\
Full 1991--2025 & Post-monsoon & rx1day mm & 35 & -0.1050 & [-0.5957, 0.3190] & 0.5027 \\
Full 1991--2025 & Post-monsoon & consecutive dry days & 35 & -0.2258 & [-0.4194, -0.0369] & 0.0840 \\
1991--2020 sensitivity & annual & precipitation sum mm & 30 & -10.6350 & [-21.3173, 0.0380] & 0.0515 \\
1991--2020 sensitivity & annual & temperature mean c & 30 & 0.0449 & [0.0022, 0.0812] & 0.0240 \\
1991--2020 sensitivity & annual & rx1day mm & 30 & -0.1611 & [-1.5650, 1.3622] & 0.5952 \\
1991--2020 sensitivity & annual & consecutive dry days & 30 & 0.0000 & [-0.3636, 0.3077] & 1.0000 \\
1996--2025 sensitivity & annual & precipitation sum mm & 30 & -3.4000 & [-8.8046, 2.5690] & 0.1909 \\
1996--2025 sensitivity & annual & temperature mean c & 30 & 0.0133 & [-0.0070, 0.0345] & 0.1779 \\
1996--2025 sensitivity & annual & rx1day mm & 30 & 0.0885 & [-1.1432, 1.5178] & 0.6087 \\
1996--2025 sensitivity & annual & consecutive dry days & 30 & 0.2727 & [-0.1864, 0.6447] & 0.1784 \\
1991--2025 sensitivity & annual & precipitation sum mm & 35 & -10.3037 & [-17.4510, -2.4821] & 0.0045 \\
1991--2025 sensitivity & annual & temperature mean c & 35 & 0.0418 & [0.0083, 0.0725] & 0.0055 \\
1991--2025 sensitivity & annual & rx1day mm & 35 & -0.6233 & [-1.5583, 0.4500] & 0.1554 \\
1991--2025 sensitivity & annual & consecutive dry days & 35 & 0.1667 & [-0.1316, 0.4334] & 0.2179 \\
\end{longtable}
\noindent\textit{Note:} Full-period p values are false-discovery-rate adjusted moving-block values. Sensitivity-row p values are unadjusted moving-block values and are shown to diagnose baseline dependence, not to multiply confirmatory tests.
\endgroup


\begin{landscape}
\begingroup
\tiny
\setlength{\tabcolsep}{3pt}
\begin{longtable}{p{0.18\textwidth}p{0.10\textwidth}p{0.16\textwidth}r r r r r r}
\caption{Complete Open-Meteo and NASA POWER comparison across temporal supports.}\label{tab:s-weather-comparison}\\
\toprule
Variable & Scale & Subset & $n$ & Pearson & Spearman & Bias & RMSE & MAE \\
\midrule
\endfirsthead
\caption[]{Complete Open-Meteo and NASA POWER comparison across temporal supports. (continued)}\\
\toprule
Variable & Scale & Subset & $n$ & Pearson & Spearman & Bias & RMSE & MAE \\
\midrule
\endhead
\midrule
\multicolumn{9}{r}{Continued on next page}\\
\endfoot
\bottomrule
\endlastfoot
precipitation sum & daily & full record & 12784 & 0.297 & 0.682 & -0.025 & 11.826 & 3.029 \\
temperature 2m mean & daily & full record & 12784 & 0.979 & 0.973 & -2.503 & 3.010 & 2.562 \\
temperature 2m max & daily & full record & 12784 & 0.954 & 0.948 & -3.535 & 4.278 & 3.667 \\
temperature 2m min & daily & full record & 12784 & 0.979 & 0.973 & -2.319 & 2.814 & 2.402 \\
relative humidity 2m mean & daily & full record & 12784 & 0.684 & 0.675 & 16.162 & 21.241 & 17.394 \\
precipitation sum & annual & full record & 35 & 0.106 & 0.083 & -9.216 & 646.739 & 431.913 \\
temperature 2m mean & annual & full record & 35 & 0.386 & 0.443 & -2.503 & 2.687 & 2.503 \\
temperature 2m max & annual & full record & 35 & 0.517 & 0.653 & -4.473 & 4.776 & 4.544 \\
precipitation sum & daily DJF & full record & 3159 & 0.254 & 0.694 & 0.585 & 12.193 & 2.063 \\
temperature 2m mean & daily DJF & full record & 3159 & 0.889 & 0.879 & -2.141 & 2.460 & 2.178 \\
temperature 2m max & daily DJF & full record & 3159 & 0.883 & 0.882 & -3.221 & 3.612 & 3.268 \\
temperature 2m min & daily DJF & full record & 3159 & 0.812 & 0.807 & -2.670 & 3.135 & 2.752 \\
relative humidity 2m mean & daily DJF & full record & 3159 & 0.686 & 0.646 & 21.260 & 24.926 & 21.466 \\
precipitation sum & daily JJ & full record & 1050 & 0.112 & 0.531 & -0.957 & 18.573 & 3.460 \\
temperature 2m mean & daily JJ & full record & 1050 & 0.754 & 0.719 & -3.537 & 4.028 & 3.590 \\
temperature 2m max & daily JJ & full record & 1050 & 0.770 & 0.761 & -4.596 & 5.201 & 4.695 \\
temperature 2m min & daily JJ & full record & 1050 & 0.696 & 0.667 & -2.950 & 3.459 & 3.007 \\
relative humidity 2m mean & daily JJ & full record & 1050 & 0.658 & 0.660 & 14.622 & 19.769 & 16.357 \\
precipitation sum & daily MAM & full record & 3220 & 0.364 & 0.585 & 0.621 & 9.378 & 2.507 \\
temperature 2m mean & daily MAM & full record & 3220 & 0.968 & 0.970 & -2.867 & 3.261 & 2.887 \\
temperature 2m max & daily MAM & full record & 3220 & 0.948 & 0.949 & -4.133 & 4.660 & 4.178 \\
temperature 2m min & daily MAM & full record & 3220 & 0.962 & 0.967 & -2.618 & 2.991 & 2.660 \\
relative humidity 2m mean & daily MAM & full record & 3220 & 0.765 & 0.760 & 14.157 & 18.143 & 15.030 \\
precipitation sum & daily Monsoon & full record & 3220 & 0.336 & 0.619 & -1.148 & 13.936 & 5.601 \\
temperature 2m mean & daily Monsoon & full record & 3220 & 0.710 & 0.675 & -2.465 & 3.196 & 2.583 \\
temperature 2m max & daily Monsoon & full record & 3220 & 0.556 & 0.503 & -3.270 & 4.529 & 3.565 \\
temperature 2m min & daily Monsoon & full record & 3220 & 0.817 & 0.777 & -1.573 & 2.149 & 1.720 \\
relative humidity 2m mean & daily Monsoon & full record & 3220 & 0.500 & 0.466 & 12.399 & 18.866 & 14.576 \\
precipitation sum & daily Post-monsoon & full record & 2135 & 0.369 & 0.579 & 0.248 & 4.709 & 1.155 \\
temperature 2m mean & daily Post-monsoon & full record & 2135 & 0.943 & 0.949 & -2.038 & 2.425 & 2.104 \\
temperature 2m max & daily Post-monsoon & full record & 2135 & 0.861 & 0.865 & -2.978 & 3.646 & 3.134 \\
temperature 2m min & daily Post-monsoon & full record & 2135 & 0.940 & 0.942 & -2.165 & 2.554 & 2.227 \\
relative humidity 2m mean & daily Post-monsoon & full record & 2135 & 0.379 & 0.324 & 18.076 & 23.594 & 19.698 \\
precipitation sum & daily & exclude either product gt200mm & 12777 & 0.425 & 0.682 & 0.135 & 8.550 & 2.846 \\
precipitation sum & daily & 1991 2019 sensitivity & 10592 & 0.467 & 0.663 & 0.612 & 8.048 & 2.715 \\
temperature 2m mean & daily & 1991 2019 sensitivity & 10592 & 0.981 & 0.976 & -2.773 & 3.213 & 2.796 \\
temperature 2m max & daily & 1991 2019 sensitivity & 10592 & 0.961 & 0.956 & -4.002 & 4.585 & 4.042 \\
temperature 2m min & daily & 1991 2019 sensitivity & 10592 & 0.979 & 0.973 & -2.421 & 2.902 & 2.489 \\
relative humidity 2m mean & daily & 1991 2019 sensitivity & 10592 & 0.727 & 0.723 & 19.163 & 22.905 & 19.412 \\
precipitation sum & annual & 1991 2019 sensitivity & 29 & 0.273 & 0.339 & 223.531 & 381.586 & 286.621 \\
temperature 2m mean & annual & 1991 2019 sensitivity & 29 & 0.502 & 0.581 & -2.773 & 2.898 & 2.773 \\
temperature 2m max & annual & 1991 2019 sensitivity & 29 & 0.707 & 0.768 & -4.961 & 5.078 & 4.961 \\
\end{longtable}
\noindent\textit{Note:} Bias is Open-Meteo minus NASA POWER. Neither product is treated as ground truth.
\endgroup
\par\medskip\noindent\textit{Panel B: Landsat-epoch agreement}\par
\begingroup
\scriptsize
\setlength{\tabcolsep}{3pt}
\begin{center}
\begin{adjustbox}{max width=\linewidth}
\begin{tabular}{p{0.34\textwidth}rrrrrr}
\toprule
Variable & $n$ & Pearson & Spearman & Bias & RMSE & MAE \\
\midrule
precipitation preceding 7d mm mean & 5 & 0.314 & 0.600 & 1.209 & 5.256 & 3.729 \\
precipitation preceding 30d mm mean & 5 & 0.371 & 0.100 & 1.798 & 16.946 & 11.962 \\
precipitation preceding 90d mm mean & 5 & -0.032 & 0.000 & -82.343 & 159.474 & 101.951 \\
precipitation preceding 180d mm mean & 5 & 0.310 & 0.700 & -276.944 & 421.897 & 294.668 \\
precipitation preceding 365d mm mean & 5 & -0.032 & 0.100 & -242.217 & 540.569 & 415.538 \\
temperature mean preceding 30d c mean & 5 & 0.758 & 0.700 & -1.687 & 1.808 & 1.687 \\
\bottomrule
\end{tabular}
\end{adjustbox}
\end{center}
\endgroup
\par\medskip\noindent\textit{Panel C: Extreme-event agreement}\par
\begingroup
\scriptsize
\setlength{\tabcolsep}{3pt}
\begin{center}
\begin{adjustbox}{max width=\linewidth}
\begin{tabular}{p{0.22\textwidth}rrrrrrr}
\toprule
Variable & OM P95 & NP P95 & OM days & NP days & Shared & Jaccard & Recall \\
\midrule
precipitation sum & 14.907 & 13.194 & 617 & 743 & 250 & 0.225 & 0.336 \\
temperature 2m max & 37.429 & 42.591 & 650 & 590 & 415 & 0.503 & 0.703 \\
\bottomrule
\end{tabular}
\end{adjustbox}
\end{center}
\endgroup
\par\medskip\noindent\textit{Panel D: Trend-slope agreement}\par
\begingroup
\scriptsize
\setlength{\tabcolsep}{3pt}
\begin{center}
\begin{adjustbox}{max width=\linewidth}
\begin{tabular}{p{0.20\textwidth}rrrp{0.11\textwidth}p{0.26\textwidth}}
\toprule
Variable & OM slope & NP slope & Same direction & Abs. difference & Interpretation \\
\midrule
precipitation sum & -10.3037 & 29.2120 & False & 39.5157 & product trend agreement; not station-validated accuracy \\
temperature 2m mean & 0.0418 & -0.0372 & False & 0.0790 & product trend agreement; not station-validated accuracy \\
temperature 2m max & 0.0333 & -0.0576 & False & 0.0909 & product trend agreement; not station-validated accuracy \\
\bottomrule
\end{tabular}
\end{adjustbox}
\end{center}
\endgroup
\end{landscape}


\section{S5: Terrain and hydrological susceptibility}

This section must document elevation, slope, roughness, terrain position, D8 flow accumulation, wetness proxy, derived drainage network and drainage distance at each analytical radius. It must report transformations and NoData rules, then preserve complete component summaries. The variables are susceptibility descriptors based on digital elevation data. They are not observed flood depth, erosion rate, blocked drainage or calibrated hazard probability.

\section{S6: Component scores and scale sensitivity}

Table~S2 must preserve every score and rank at all three analytical supports, including ties and the component identifiers used in the vector inventory. The main text reports the fixed 500~m sequence, while the supplementary table allows the reader to inspect every reversal. Rank correlations, top-three overlap, top-five overlap and maximum movement must be reported with the exact paired scales.


\section{S7: Redundancy, ablation and rank uncertainty}

Tables~S5--S6 must include the complete indicator-correlation matrix, all factor and domain removals, 50,000 Dirichlet weight draws and 2,000 spatial-block rank draws. Sampling-based uncertainty and sensitivity address different questions and must remain separately labelled \citep{helton2006sensitivity}. The fixed Giri score, the weight-perturbation distribution and the spatial-bootstrap Bhallar/Giri overlap must not be collapsed into a single winning component.

\begin{landscape}
\begingroup
\scriptsize
\setlength{\tabcolsep}{3pt}
\begin{longtable}{p{0.22\textwidth}rrrp{0.25\textwidth}r}
\caption{Factor- and domain-ablation sensitivity.}\label{tab:s-ablation}\\
\toprule
Scenario & Spearman & Kendall & Top 5 & Rank 1 component & Max shift \\
\midrule
\endfirsthead
\caption[]{Factor- and domain-ablation sensitivity. (continued)}\\
\toprule
Scenario & Spearman & Kendall & Top 5 & Rank 1 component & Max shift \\
\midrule
\endhead
\midrule
\multicolumn{6}{r}{Continued on next page}\\
\endfoot
\bottomrule
\endlastfoot
Full integrated context & 1.000 & 1.000 & 5 & Giri complex of monuments & 0.0 \\
Primary local priority & 1.000 & 1.000 & 5 & Giri complex of monuments & 0.0 \\
Landscape only & 0.325 & 0.185 & 1 & Sirkap (fortified city) & 11.0 \\
Terrain only & 0.830 & 0.687 & 4 & Buddhist remains around Bhallar stupa & 7.0 \\
Climate only &  &  & 0 & All components tied &  \\
Leave climate out & 1.000 & 1.000 & 5 & Giri complex of monuments & 0.0 \\
Leave landscape out & 0.830 & 0.687 & 4 & Buddhist remains around Bhallar stupa & 7.0 \\
Leave terrain out & 0.325 & 0.185 & 1 & Sirkap (fortified city) & 11.0 \\
Leave NDVI out & 0.992 & 0.952 & 5 & Giri complex of monuments & 1.0 \\
Leave MNDWI out & -0.280 & -0.214 & 0 & Buddhist remains around Bhallar stupa & 13.0 \\
Leave NDBI out & 0.990 & 0.952 & 4 & Giri complex of monuments & 2.0 \\
Leave slope out & 0.611 & 0.435 & 3 & Sirsukh (fortified ruined city) & 8.0 \\
Leave terrain wetness out & 0.892 & 0.748 & 4 & Giri Mosque and tombs & 4.0 \\
Leave drainage proximity out & 0.460 & 0.311 & 2 & Sirsukh (fortified ruined city) & 8.5 \\
\end{longtable}
\noindent\textit{Note:} Panel A compares each scenario with the full integrated ordering. Panels B and C report all component ranks.
\endgroup
\par\medskip\noindent\textit{Panel B: Domain-ablation component ranks}\par
\begingroup
\tiny
\setlength{\tabcolsep}{3pt}
\begin{center}
\begin{adjustbox}{max width=\linewidth}
\begin{tabular}{p{0.08\textwidth}p{0.25\textwidth}rrrrrrr}
\toprule
ID & Component & Full & Local & Landscape & Terrain & Climate & No climate & No landscape \\
\midrule
139-001 & Khanpur Cave & 14.0 & 14.0 & 14.0 & 12.0 & 9.0 & 14.0 & 12.0 \\
139-003 & Bhir Mound & 17.0 & 17.0 & 16.0 & 14.0 & 9.0 & 17.0 & 14.0 \\
139-004 & Sirkap (fortified city) & 12.0 & 12.0 & 1.0 & 15.0 & 9.0 & 12.0 & 15.0 \\
139-005 & Sirsukh (fortified ruined city) & 9.0 & 9.0 & 3.0 & 12.0 & 9.0 & 9.0 & 12.0 \\
139-006 & Dharmarajika stupa and monastery & 5.0 & 5.0 & 5.0 & 8.0 & 9.0 & 5.0 & 8.0 \\
139-007 & Khader Mohra (Akhuri) & 6.0 & 6.0 & 8.0 & 6.0 & 9.0 & 6.0 & 6.0 \\
139-008 & Kalawan group of buildings & 7.0 & 7.0 & 2.0 & 9.0 & 9.0 & 7.0 & 9.0 \\
139-009 & Giri complex of monuments & 1.0 & 1.0 & 9.0 & 2.0 & 9.0 & 1.0 & 2.0 \\
139-010 & Kunala stupa and monastery & 10.5 & 10.5 & 6.0 & 10.0 & 9.0 & 10.5 & 10.0 \\
139-011 & Jandial complex & 16.0 & 16.0 & 17.0 & 12.0 & 9.0 & 16.0 & 12.0 \\
139-012 & Lalchak and Badalpur Buddhist stuppa & 15.0 & 15.0 & 12.0 & 17.0 & 9.0 & 15.0 & 17.0 \\
139-013 & Mohra Moradu stupa and monastery & 4.0 & 4.0 & 10.0 & 5.0 & 9.0 & 4.0 & 5.0 \\
139-014 & Pippala stupa and monastery & 10.5 & 10.5 & 13.0 & 7.0 & 9.0 & 10.5 & 7.0 \\
139-015 & Jaulian stupa and monastery & 3.0 & 3.0 & 7.0 & 4.0 & 9.0 & 3.0 & 4.0 \\
139-016 & Lalchak mounds & 13.0 & 13.0 & 4.0 & 16.0 & 9.0 & 13.0 & 16.0 \\
139-017 & Buddhist remains around Bhallar stupa & 8.0 & 8.0 & 15.0 & 1.0 & 9.0 & 8.0 & 1.0 \\
139-018 & Giri Mosque and tombs & 2.0 & 2.0 & 11.0 & 3.0 & 9.0 & 2.0 & 3.0 \\
\bottomrule
\end{tabular}
\end{adjustbox}
\end{center}
\endgroup
\par\medskip\noindent\textit{Panel C: Factor-ablation component ranks}\par
\begingroup
\tiny
\setlength{\tabcolsep}{3pt}
\begin{center}
\begin{adjustbox}{max width=\linewidth}
\begin{tabular}{p{0.08\textwidth}p{0.25\textwidth}rrrrrrr}
\toprule
ID & Component & No terrain & No NDVI & No MNDWI & No NDBI & No slope & No wetness & No drainage \\
\midrule
139-001 & Khanpur Cave & 14.0 & 14.0 & 8.0 & 14.0 & 16.0 & 11.0 & 15.0 \\
139-003 & Bhir Mound & 16.0 & 17.0 & 4.0 & 17.0 & 17.0 & 17.0 & 16.0 \\
139-004 & Sirkap (fortified city) & 1.0 & 12.0 & 3.0 & 12.0 & 10.0 & 12.0 & 7.0 \\
139-005 & Sirsukh (fortified ruined city) & 3.0 & 9.0 & 14.0 & 11.0 & 1.0 & 13.0 & 1.0 \\
139-006 & Dharmarajika stupa and monastery & 5.0 & 4.0 & 7.0 & 6.0 & 2.0 & 9.0 & 3.0 \\
139-007 & Khader Mohra (Akhuri) & 8.0 & 6.0 & 11.0 & 5.0 & 11.0 & 4.0 & 13.0 \\
139-008 & Kalawan group of buildings & 2.0 & 8.0 & 10.0 & 7.0 & 6.0 & 8.0 & 9.0 \\
139-009 & Giri complex of monuments & 9.0 & 1.0 & 13.0 & 1.0 & 4.0 & 2.0 & 8.0 \\
139-010 & Kunala stupa and monastery & 6.0 & 10.0 & 6.0 & 10.0 & 14.0 & 6.5 & 2.0 \\
139-011 & Jandial complex & 17.0 & 16.0 & 5.0 & 16.0 & 15.0 & 15.0 & 17.0 \\
139-012 & Lalchak and Badalpur Buddhist stuppa & 12.0 & 15.0 & 17.0 & 15.0 & 13.0 & 16.0 & 14.0 \\
139-013 & Mohra Moradu stupa and monastery & 10.0 & 5.0 & 9.0 & 4.0 & 12.0 & 5.0 & 10.0 \\
139-014 & Pippala stupa and monastery & 13.0 & 11.0 & 2.0 & 9.0 & 8.0 & 6.5 & 12.0 \\
139-015 & Jaulian stupa and monastery & 7.0 & 2.0 & 12.0 & 3.0 & 5.0 & 3.0 & 11.0 \\
139-016 & Lalchak mounds & 4.0 & 13.0 & 16.0 & 13.0 & 7.0 & 14.0 & 5.5 \\
139-017 & Buddhist remains around Bhallar stupa & 15.0 & 7.0 & 1.0 & 8.0 & 3.0 & 10.0 & 5.5 \\
139-018 & Giri Mosque and tombs & 11.0 & 3.0 & 15.0 & 2.0 & 9.0 & 1.0 & 4.0 \\
\bottomrule
\end{tabular}
\end{adjustbox}
\end{center}
\endgroup
\end{landscape}


\begin{landscape}
\begingroup
\scriptsize
\setlength{\tabcolsep}{3pt}
\begin{longtable}{p{0.09\textwidth}p{0.30\textwidth}rrp{0.14\textwidth}rr}
\caption{Complete domain-weight and spatial-block rank distributions.}\label{tab:s-rank-uncertainty}\\
\toprule
ID & Component & Fixed R & Median R & 95\% R interval & P(top 3) & P(top 5) \\
\midrule
\endfirsthead
\caption[]{Complete domain-weight and spatial-block rank distributions. (continued)}\\
\toprule
ID & Component & Fixed R & Median R & 95\% R interval & P(top 3) & P(top 5) \\
\midrule
\endhead
\midrule
\multicolumn{7}{r}{Continued on next page}\\
\endfoot
\bottomrule
\endlastfoot
139-009 & Giri complex of monuments & 1.0 & 1.0 & [1.0, 9.0] & 0.774 & 0.832 \\
139-018 & Giri Mosque and tombs & 2.0 & 3.0 & [2.0, 11.0] & 0.679 & 0.730 \\
139-015 & Jaulian stupa and monastery & 3.0 & 4.0 & [2.0, 7.0] & 0.440 & 0.834 \\
139-013 & Mohra Moradu stupa and monastery & 4.0 & 5.0 & [4.0, 11.0] & 0.000 & 0.614 \\
139-006 & Dharmarajika stupa and monastery & 5.0 & 5.0 & [2.0, 8.0] & 0.219 & 0.512 \\
139-007 & Khader Mohra (Akhuri) & 6.0 & 6.0 & [5.0, 9.0] & 0.000 & 0.045 \\
139-008 & Kalawan group of buildings & 7.0 & 7.0 & [2.0, 9.0] & 0.254 & 0.386 \\
139-017 & Buddhist remains around Bhallar stupa & 8.0 & 8.0 & [1.0, 15.0] & 0.305 & 0.443 \\
139-005 & Sirsukh (fortified ruined city) & 9.0 & 9.0 & [3.0, 11.0] & 0.102 & 0.202 \\
139-014 & Pippala stupa and monastery & 10.0 & 10.0 & [7.0, 13.0] & 0.000 & 0.000 \\
139-010 & Kunala stupa and monastery & 11.0 & 10.0 & [6.0, 10.0] & 0.000 & 0.009 \\
139-004 & Sirkap (fortified city) & 12.0 & 12.0 & [1.0, 15.0] & 0.226 & 0.270 \\
139-016 & Lalchak mounds & 13.0 & 13.0 & [4.0, 16.0] & 0.000 & 0.123 \\
139-001 & Khanpur Cave & 14.0 & 14.0 & [12.0, 15.0] & 0.000 & 0.000 \\
139-012 & Lalchak and Badalpur Buddhist stuppa & 15.0 & 15.0 & [12.0, 17.0] & 0.000 & 0.000 \\
139-011 & Jandial complex & 16.0 & 16.0 & [13.0, 17.0] & 0.000 & 0.000 \\
139-003 & Bhir Mound & 17.0 & 16.0 & [14.0, 17.0] & 0.000 & 0.000 \\
\end{longtable}
\noindent\textit{Note:} Panel A uses 50,000 Dirichlet domain-weight draws. Panel B uses 2,000 resamples of 150 m spectral blocks with terrain held fixed.
\endgroup
\par\medskip\noindent\textit{Panel B: Spatial-block rank distribution}\par
\begingroup
\scriptsize
\setlength{\tabcolsep}{3pt}
\begin{center}
\begin{adjustbox}{max width=\linewidth}
\begin{tabular}{p{0.09\textwidth}p{0.27\textwidth}rp{0.14\textwidth}rp{0.14\textwidth}rr}
\toprule
ID & Component & Median S & 95\% S interval & Median R & 95\% R interval & P(top 3) & P(top 5) \\
\midrule
139-017 & Buddhist remains around Bhallar stupa & 0.598 & [0.547, 0.627] & 1.0 & [1.0, 5.0] & 0.911 & 0.989 \\
139-009 & Giri complex of monuments & 0.603 & [0.581, 0.603] & 2.0 & [1.0, 3.0] & 0.986 & 1.000 \\
139-018 & Giri Mosque and tombs & 0.564 & [0.564, 0.586] & 3.0 & [2.0, 5.0] & 0.615 & 0.997 \\
139-015 & Jaulian stupa and monastery & 0.559 & [0.515, 0.581] & 4.0 & [2.0, 8.5] & 0.294 & 0.838 \\
139-006 & Dharmarajika stupa and monastery & 0.532 & [0.488, 0.576] & 6.0 & [3.0, 11.0] & 0.038 & 0.376 \\
139-013 & Mohra Moradu stupa and monastery & 0.510 & [0.436, 0.583] & 8.0 & [3.0, 15.0] & 0.051 & 0.218 \\
139-014 & Pippala stupa and monastery & 0.510 & [0.473, 0.569] & 8.0 & [3.5, 12.0] & 0.025 & 0.148 \\
139-010 & Kunala stupa and monastery & 0.510 & [0.466, 0.554] & 8.0 & [4.5, 12.0] & 0.005 & 0.089 \\
139-007 & Khader Mohra (Akhuri) & 0.500 & [0.434, 0.588] & 9.0 & [3.0, 15.0] & 0.046 & 0.175 \\
139-008 & Kalawan group of buildings & 0.500 & [0.448, 0.551] & 9.0 & [4.5, 14.0] & 0.004 & 0.086 \\
139-005 & Sirsukh (fortified ruined city) & 0.500 & [0.456, 0.530] & 9.0 & [5.5, 13.0] & 0.001 & 0.023 \\
139-004 & Sirkap (fortified city) & 0.490 & [0.461, 0.534] & 10.0 & [6.0, 13.0] & 0.001 & 0.018 \\
139-016 & Lalchak mounds & 0.446 & [0.417, 0.475] & 14.0 & [11.0, 16.0] & 0.000 & 0.000 \\
139-011 & Jandial complex & 0.438 & [0.404, 0.471] & 14.0 & [12.0, 16.0] & 0.000 & 0.000 \\
139-001 & Khanpur Cave & 0.441 & [0.404, 0.507] & 14.0 & [8.5, 16.0] & 0.001 & 0.002 \\
139-003 & Bhir Mound & 0.426 & [0.390, 0.449] & 16.0 & [13.0, 16.0] & 0.000 & 0.000 \\
139-012 & Lalchak and Badalpur Buddhist stuppa & 0.363 & [0.333, 0.407] & 17.0 & [16.0, 17.0] & 0.000 & 0.000 \\
\bottomrule
\end{tabular}
\end{adjustbox}
\end{center}
\endgroup
\end{landscape}


\section{S8: Geographic proxy-label experiment}

Table~S7 must preserve development, buffer and test sample counts, group definitions, tuning records, all model metrics, class-wise confusion, calibration and block-bootstrap intervals. Spatial validation must be aligned with the geographic transfer claim \citep{hijmans2012cv,meyer2018target,meyer2019selection,meyer2021aoa,roberts2017cv}. Macro-F1, multiclass area under the receiver-operating-characteristic curve and probability scores describe complementary properties \citep{saito2015precision,hand2001auc,gneiting2007scoring}. None can overcome the lack of independent condition labels. The accuracy claim is restricted to WorldCover-derived proxy agreement, consistent with established land-cover accuracy and disagreement literature \citep{congalton1991accuracy,foody2002accuracy,olofsson2014accuracy,pontius2011kappa}.

\begingroup
\tiny
\setlength{\tabcolsep}{3pt}
\begin{longtable}{p{0.22\textwidth}r r p{0.13\textwidth}rrrrr}
\caption{Complete geographically buffered proxy-model metrics, calibration and confusion records.}\label{tab:s-proxy-model}\\
\toprule
Model & Inner F1 & Test F1 & 95\% interval & Bal. acc. & AUROC & AUPRC & Log loss & ECE \\
\midrule
\endfirsthead
\caption[]{Complete geographically buffered proxy-model metrics, calibration and confusion records. (continued)}\\
\toprule
Model & Inner F1 & Test F1 & 95\% interval & Bal. acc. & AUROC & AUPRC & Log loss & ECE \\
\midrule
\endhead
\midrule
\multicolumn{9}{r}{Continued on next page}\\
\endfoot
\bottomrule
\endlastfoot
MLP sensitivity & 0.871 & 0.831 & [0.700, 0.860] & 0.810 & 0.977 & 0.910 & 0.459 & 0.014 \\
Histogram gradient boosting & 0.863 & 0.821 & [0.693, 0.850] & 0.799 & 0.974 & 0.909 & 0.483 & 0.041 \\
Extra trees & 0.861 & 0.826 & [0.685, 0.838] & 0.808 & 0.973 & 0.904 & 0.498 & 0.017 \\
Random forest & 0.860 & 0.829 & [0.689, 0.842] & 0.810 & 0.973 & 0.906 & 0.500 & 0.016 \\
Regularised multinomial logistic & 0.844 & 0.842 & [0.676, 0.860] & 0.835 & 0.971 & 0.900 & 0.527 & 0.035 \\
\end{longtable}
\noindent\textit{Note:} All outcomes measure agreement with WorldCover-derived proxy land-cover labels. They do not measure independent heritage-condition accuracy.
\endgroup
\par\medskip\noindent\textit{Panel B: Proxy-class sample distribution}\par
\begingroup
\scriptsize
\setlength{\tabcolsep}{3pt}
\begin{center}
\begin{adjustbox}{max width=\linewidth}
\begin{tabular}{lrrr}
\toprule
Class & Development & Buffer-excluded & Test \\
\midrule
1 & 1885 & 1228 & 887 \\
2 & 2604 & 666 & 730 \\
3 & 2182 & 881 & 937 \\
4 & 2442 & 873 & 685 \\
5 & 2340 & 801 & 859 \\
6 & 641 & 18 & 90 \\
\bottomrule
\end{tabular}
\end{adjustbox}
\end{center}
\endgroup
\par\medskip\noindent\textit{Panel C: Primary-model test confusion matrix}\par
\begingroup
\scriptsize
\setlength{\tabcolsep}{3pt}
\begin{center}
\begin{adjustbox}{max width=\linewidth}
\begin{tabular}{lrrrrrr}
\toprule
Reference & Pred. 1 & Pred. 2 & Pred. 3 & Pred. 4 & Pred. 5 & Pred. 6 \\
\midrule
reference\_1 & 729 & 63 & 93 & 0 & 2 & 0 \\
reference\_2 & 56 & 663 & 8 & 0 & 2 & 1 \\
reference\_3 & 39 & 45 & 772 & 21 & 60 & 0 \\
reference\_4 & 1 & 0 & 26 & 569 & 89 & 0 \\
reference\_5 & 6 & 21 & 124 & 12 & 696 & 0 \\
reference\_6 & 1 & 29 & 0 & 0 & 0 & 60 \\
\bottomrule
\end{tabular}
\end{adjustbox}
\end{center}
\endgroup
\par\medskip\noindent\textit{Panel D: Calibration for Regularised multinomial logistic}\par
\begingroup
\scriptsize
\setlength{\tabcolsep}{3pt}
\begin{center}
\begin{adjustbox}{max width=\linewidth}
\begin{tabular}{rrrrr}
\toprule
Bin & Count & Accuracy & Confidence & ECE \\
\midrule
0.0--0.1 & 0 &  &  & 0.035 \\
0.1--0.2 & 0 &  &  & 0.035 \\
0.2--0.3 & 1 & 0.000 & 0.297 & 0.035 \\
0.3--0.4 & 31 & 0.355 & 0.372 & 0.035 \\
0.4--0.5 & 215 & 0.447 & 0.463 & 0.035 \\
0.5--0.6 & 579 & 0.582 & 0.551 & 0.035 \\
0.6--0.7 & 589 & 0.705 & 0.651 & 0.035 \\
0.7--0.8 & 613 & 0.811 & 0.752 & 0.035 \\
0.8--0.9 & 717 & 0.914 & 0.850 & 0.035 \\
0.9--1.0 & 1443 & 0.971 & 0.965 & 0.035 \\
\bottomrule
\end{tabular}
\end{adjustbox}
\end{center}
\endgroup
\par\medskip\noindent\textit{Panel D: Calibration for Random forest}\par
\begingroup
\scriptsize
\setlength{\tabcolsep}{3pt}
\begin{center}
\begin{adjustbox}{max width=\linewidth}
\begin{tabular}{rrrrr}
\toprule
Bin & Count & Accuracy & Confidence & ECE \\
\midrule
0.0--0.1 & 0 &  &  & 0.016 \\
0.1--0.2 & 0 &  &  & 0.016 \\
0.2--0.3 & 0 &  &  & 0.016 \\
0.3--0.4 & 32 & 0.438 & 0.367 & 0.016 \\
0.4--0.5 & 159 & 0.478 & 0.460 & 0.016 \\
0.5--0.6 & 402 & 0.585 & 0.552 & 0.016 \\
0.6--0.7 & 386 & 0.679 & 0.650 & 0.016 \\
0.7--0.8 & 548 & 0.748 & 0.754 & 0.016 \\
0.8--0.9 & 851 & 0.826 & 0.854 & 0.016 \\
0.9--1.0 & 1810 & 0.962 & 0.967 & 0.016 \\
\bottomrule
\end{tabular}
\end{adjustbox}
\end{center}
\endgroup
\par\medskip\noindent\textit{Panel D: Calibration for Extra trees}\par
\begingroup
\scriptsize
\setlength{\tabcolsep}{3pt}
\begin{center}
\begin{adjustbox}{max width=\linewidth}
\begin{tabular}{rrrrr}
\toprule
Bin & Count & Accuracy & Confidence & ECE \\
\midrule
0.0--0.1 & 0 &  &  & 0.017 \\
0.1--0.2 & 0 &  &  & 0.017 \\
0.2--0.3 & 0 &  &  & 0.017 \\
0.3--0.4 & 46 & 0.457 & 0.376 & 0.017 \\
0.4--0.5 & 190 & 0.511 & 0.462 & 0.017 \\
0.5--0.6 & 456 & 0.586 & 0.550 & 0.017 \\
0.6--0.7 & 451 & 0.687 & 0.650 & 0.017 \\
0.7--0.8 & 621 & 0.778 & 0.751 & 0.017 \\
0.8--0.9 & 808 & 0.855 & 0.850 & 0.017 \\
0.9--1.0 & 1616 & 0.970 & 0.967 & 0.017 \\
\bottomrule
\end{tabular}
\end{adjustbox}
\end{center}
\endgroup
\par\medskip\noindent\textit{Panel D: Calibration for Histogram gradient boosting}\par
\begingroup
\scriptsize
\setlength{\tabcolsep}{3pt}
\begin{center}
\begin{adjustbox}{max width=\linewidth}
\begin{tabular}{rrrrr}
\toprule
Bin & Count & Accuracy & Confidence & ECE \\
\midrule
0.0--0.1 & 0 &  &  & 0.041 \\
0.1--0.2 & 0 &  &  & 0.041 \\
0.2--0.3 & 0 &  &  & 0.041 \\
0.3--0.4 & 16 & 0.562 & 0.376 & 0.041 \\
0.4--0.5 & 96 & 0.417 & 0.464 & 0.041 \\
0.5--0.6 & 297 & 0.485 & 0.549 & 0.041 \\
0.6--0.7 & 308 & 0.617 & 0.654 & 0.041 \\
0.7--0.8 & 411 & 0.703 & 0.754 & 0.041 \\
0.8--0.9 & 633 & 0.798 & 0.855 & 0.041 \\
0.9--1.0 & 2427 & 0.941 & 0.972 & 0.041 \\
\bottomrule
\end{tabular}
\end{adjustbox}
\end{center}
\endgroup
\par\medskip\noindent\textit{Panel D: Calibration for MLP sensitivity}\par
\begingroup
\scriptsize
\setlength{\tabcolsep}{3pt}
\begin{center}
\begin{adjustbox}{max width=\linewidth}
\begin{tabular}{rrrrr}
\toprule
Bin & Count & Accuracy & Confidence & ECE \\
\midrule
0.0--0.1 & 0 &  &  & 0.014 \\
0.1--0.2 & 0 &  &  & 0.014 \\
0.2--0.3 & 0 &  &  & 0.014 \\
0.3--0.4 & 27 & 0.333 & 0.371 & 0.014 \\
0.4--0.5 & 121 & 0.430 & 0.468 & 0.014 \\
0.5--0.6 & 334 & 0.530 & 0.552 & 0.014 \\
0.6--0.7 & 365 & 0.658 & 0.651 & 0.014 \\
0.7--0.8 & 458 & 0.729 & 0.752 & 0.014 \\
0.8--0.9 & 735 & 0.848 & 0.854 & 0.014 \\
0.9--1.0 & 2148 & 0.956 & 0.970 & 0.014 \\
\bottomrule
\end{tabular}
\end{adjustbox}
\end{center}
\endgroup


\section{S9: Heritage interpretation and field use}

The supplementary interpretation record must distinguish screening from complete risk assessment, and it must keep landscape change separate from material condition. Heritage Earth-observation studies support repeated landscape surveillance and change detection \citep{agapiou2017petabyte,agapiou2021dynamic,tapete2017sar}. Climate-vulnerability and adaptation research further shows that usable conservation decisions require vulnerability, capacity, local knowledge and field evidence beyond remote forcing layers \citep{sesana2020vulnerability,crowley2022risk,daly2022adaptation}. These sources support the paper's conservative field-inspection endpoint; they do not supply missing condition labels for Taxila.

\section{S10: Reproducibility and experiment disposition}

The final supplementary archive must include or cross-reference the source manifest, scene and weather inventories, data dictionary, configurations, random seed, dependency record, execution logs, checksums, claim-evidence audit, equation-code map, display-source map, discrepancy record, result-disposition matrix and limitations register. Fourteen experiments were completed, the weather-spectral association experiment remained exploratory, and the ancillary-hazard experiment was not executed. Earthquake, flood-event, landslide, quarry and population layers must not appear as completed analyses.

\section{Caption and reporting checklist}

Every final display must define abbreviations, units, denominator, uncertainty and analysis scale. Maps must state the coordinate reference system and identify whether they represent pressure, susceptibility or priority. Captions must explain NoData and source provenance. Table notes must distinguish fixed scores, ranks, percentile intervals and probabilities. No caption may use damage, condition accuracy or risk unless the displayed analysis independently supports that term. Final insertion, numbering and cross-referencing of the verified files are Phase~8 assembly tasks.
